% Generated by scripts/assemble.py; corrections belong in source/corrections.json.
\examhead{157}{华南理工大学 2026 年数学分析试卷}{150}
% p151-b04; source PDF page 151
\par\Needspace{5\baselineskip}
\blockanchor{p151-b04}
\majorquestion{一}证明题. 共114分.\par
% p151-b05; source PDF page 151
\blockanchor{p151-b05}
\question{1}设  $x_{0}=2,\quad x_{n+1}=\frac{1}{2}\left(x_{n}+\frac{3}{x_{n}}\right)\quad(n=0,1,\cdots)$ ，证明： $\lim_{n\to\infty}x_{n}=\sqrt{3}$ .\par
% p151-b06; source PDF page 151
\blockanchor{p151-b06}
\question{2}$f(x) = x^{2}\sin x$ 在 $0 \leq x < +\infty$ 上是否一致连续？给出证明.\par
% p151-b07; source PDF page 151
\par\Needspace{6\baselineskip}
\blockanchor{p151-b07}
\question{3}$f$ 在 $[0,2]$ 上二阶连续可微, 且 $|f''(x)| \leq M$. 若 $f$ 在 $(0,2)$ 内能取到最大值, 证明:\par
\nopagebreak[4]
% p151-b08; source PDF page 151
\blockanchor{p151-b08}
\displayformula{\left| f ^ {\prime} (0) + f ^ {\prime} (2) \right| \leq 2 M.}
% p151-b09; source PDF page 151
\blockanchor{p151-b09}
\question{4}已知 $f(x)$ 在 $[A, B]$ 上黎曼可积, 证明: $\lim_{h \to 0} \int_{a}^{b} |f(x + h) - f(x)| \mathrm{d}x = 0$, 其中 $A < a < b < B$.\par
% p151-b10; source PDF page 151
\blockanchor{p151-b10}
\question{5}设 $f(t)$ 在 $(- \infty, + \infty)$ 上连续且绝对可积. 证明: $g(w) = \int_{-\infty}^{+\infty} f(t) \cos(wt) \, dt$ 在 $w \in (-\infty, + \infty)$ 上有界且一致连续.\par
% p151-b11; source PDF page 151
\blockanchor{p151-b11}
\question{6}判断 $\sum_{n=2}^{\infty} \ln \left(1 + \frac{(-1)^n}{n^\alpha}\right) (\alpha > 0)$ 的绝对收敛性, 条件收敛性, 发散性.\par
% p151-b12; source PDF page 151
\blockanchor{p151-b12}
\question{7}解答如下问题:\par
% p151-b13; source PDF page 151
\blockanchor{p151-b13}
(1) 求 $|\sin x|$ 在 $(- \infty, + \infty)$ 上的傅里叶级数, 该级数是否一致收敛?\par
% p151-b14; source PDF page 151
\blockanchor{p151-b14}
(2) 若 $f(x)$ 在 $[a, b]$ 上黎曼可积, 证明: $\lim_{\lambda \to +\infty} \int_{a}^{b} f(x) |\sin \lambda x| \mathrm{d}x = \frac{2}{\pi} \int_{a}^{b} f(x) \mathrm{d}x.$\par
% p151-b15; source PDF page 151
\blockanchor{p151-b15}
\question{8}记区域 $D: 0 \leq x \leq 1, 0 \leq y \leq 1$, 证明: $e^{xy} + \ln(1 + xy) \geq 1 + xy$, 并说明取等条件.\par
% p151-b16; source PDF page 151
\par\Needspace{6\baselineskip}
\blockanchor{p151-b16}
\question{9}已知\par
\nopagebreak[4]
% p151-b17; source PDF page 151
\blockanchor{p151-b17}
\displayformula{f (x, y) = \left\{ \begin{array}{l l} { \frac {x ^ {3}}{x ^ {2} + y ^ {2}},} & {x ^ {2} + y ^ {2} \neq 0} \\ {0,} & {x ^ {2} + y ^ {2} \neq 0 (\text {应该是} x ^ {2} + y ^ {2} = 0)} \end{array} \right.}
% p151-b18; source PDF page 151
\blockanchor{p151-b18}
(1) 若 $x = x(t), y = y(t)$ 是过原点的任意可微曲线 (即 $x^{2}(0) + y^{2}(0) = 0$, 且 $(x'(t))^{2} + (y'(t))^{2} > 0$), 证明: $f(x(t), y(t))$ 可微.\par
% p151-b19; source PDF page 151
\blockanchor{p151-b19}
(2) 求证 $f(x, y)$ 在 $(0, 0)$ 处不可微.\par
% p151-b20; source PDF page 151
\par\Needspace{5\baselineskip}
\blockanchor{p151-b20}
\majorquestion{二}计算题. 共 36 分.\par
% p151-b21; source PDF page 151
\blockanchor{p151-b21}
\question{1}求 $\lim_{x\to 0}\frac{(5 + x)^x - 5^x}{x^2}.$\par
% p151-b22; source PDF page 151
\par\Needspace{6\baselineskip}
\blockanchor{p151-b22}
\question{2}求\par
\nopagebreak[4]
% p151-b23; source PDF page 151
\blockanchor{p151-b23}
\displayformula{I = \int_ {L} (x ^ {2} + z ^ {2}) \mathrm{d} y + (y ^ {2} + x ^ {2}) \mathrm{d} z.}
% p151-b24; source PDF page 151
\blockanchor{p151-b24}
$L$ 是 $x^{2} + y^{2} + z^{2} = 2ax$ 与 $y^{2} + x^{2} = 2bx$ 的交线, 方向是经由 $y > 0$ 的方向再回到原点 $(a > b)$.\par
% p151-b25; source PDF page 151
\par\Needspace{6\baselineskip}
\blockanchor{p151-b25}
\question{3}(可能有误)已知 $f(x, y, z) = x^2 + y^2 - z$, $\Sigma$ 是一光滑封闭曲面, $(x_0, y_0, z_0)$ 是单位球体上的一点, $n$ 是 $\Sigma$ 的单位外法向量, 计算\par
\nopagebreak[4]
% p151-b26; source PDF page 151
\blockanchor{p151-b26}
\displayformula{I = \iint_ {\Sigma} \left(f \frac {\cos (\widehat {r , n})}{r ^ {2}} + \frac {1}{r} \frac {\partial f}{\partial n}\right) d S.}
% p151-b27; source PDF page 151
\blockanchor{p151-b27}
其中  $\boldsymbol{r}=(x-x_{0},y-y_{0},z-z_{0}),\boldsymbol{r}=|\boldsymbol{r}|$ . (I 应该无误, 其余  $\Sigma$ , n 等有关叙述可能不太准则)\par
